% ============================================================================
% Fichier  : partie3/P03_C12_post_exploitation.tex
% Titre    : Post-Exploitation, Mouvement Latéral, Persistance
% Auteur   : MINKA MI NGUIDJOÏ Thierry Emmanuel
% Version  : 1.0 -- novembre 2026
% ============================================================================
% Sources : MITRE ATT&CK (tactiques Lateral Movement, Persistence, Collection,
% Exfiltration) ; journaux de sécurité Windows (événements 4624, 4672, 4698,
% 4720, 7045). Opérations GAROUA et LOGONE : cas fictifs.
% ============================================================================

\chapter{Post-Exploitation, Mouvement Latéral, Persistance}
\label{chap:post}
\chaptermark{Post-exploitation}

\epigraph{%
    « Le testeur laisse le système tel qu'il l'a trouvé~;\\
    l'attaquant y laisse une porte. »%
}{Principe du cours}

\niveauChapitre{Expert}{Chapitres~\ref{chap:ad} et~\ref{chap:privileges}.}

\begin{boxobjectifs}
\begin{itemize}
    \item Décrire les objectifs de la post-exploitation~: comprendre,
          s'étendre, durer, atteindre les données.
    \item Expliquer les techniques de mouvement latéral et leurs traces.
    \item Reconnaître les mécanismes de persistance.
    \item Fixer les limites du testeur~: démontrer sans installer.
    \item Relier chaque technique aux événements qui la détectent.
\end{itemize}
\end{boxobjectifs}

\bigskip

% ============================================================================
\section*{Question directrice : que fait l'attaquant après être devenu administrateur, et que doit faire le testeur~?}
% ============================================================================

En mars, l'attaquant devenu administrateur a créé un compte, installé un
service, copié 40~Go et détruit les sauvegardes. En octobre, Mme~NGO~MBOCK
devenue administratrice s'est arrêtée dans les cinq minutes. Ce chapitre
décrit ce qu'un attaquant fait après l'élévation, pour que le défenseur le
voie, et ce que le testeur s'interdit.

\medskip
\textbf{Plan du chapitre.}
\begin{enumerate}
    \item Les objectifs de la post-exploitation (section~\ref{sec:c12-objectifs}).
    \item Le mouvement latéral (section~\ref{sec:c12-lateral}).
    \item La persistance (section~\ref{sec:c12-persistance}).
    \item Les limites du testeur (section~\ref{sec:c12-limites}).
    \item Ce que voient les journaux (section~\ref{sec:c12-journaux}).
\end{enumerate}

\separateur

% ============================================================================
\section{Les Objectifs de la Post-Exploitation}
\label{sec:c12-objectifs}
% ============================================================================

\begin{tableaucours}{L{3.4cm}L{4.4cm}L{4.4cm}}{%
    \tableauHeader{Objectif} &
    \tableauHeader{Tactique ATT\&CK} &
    \tableauHeader{Dans l'Opération GAROUA}
}
Comprendre l'environnement & Découverte & Balayage du réseau, interrogation
de l'annuaire \\
S'étendre & Mouvement latéral & Connexions avec le mot de passe local
partagé \\
Durer & Persistance & Compte svc\_maj2, service WinUpd32 \\
Atteindre les données & Collecte, exfiltration & Copie de 40~Go \\
Agir & Impact & Suppression des sauvegardes, chiffrement \\
\end{tableaucours}
\vspace{-0.8em}
\begin{center}{\small\itshape Tableau~12.1. Les objectifs de la post-exploitation~\cite{mitreattack}}\end{center}

% ============================================================================
\section{Le Mouvement Latéral}
\label{sec:c12-lateral}
% ============================================================================

\begin{boxdef}
\termecle{Mouvement latéral}~: ensemble des techniques par lesquelles un
attaquant, présent sur une machine, prend le contrôle d'autres machines du
réseau, le plus souvent avec des identifiants valides.
\end{boxdef}

\begin{tableaucours}{L{4.0cm}L{4.4cm}L{3.8cm}}{%
    \tableauHeader{Technique} &
    \tableauHeader{Principe} &
    \tableauHeader{Trace principale}
}
Services d'administration à distance & Exécution d'une commande par
création de service distant & Événement 7045 (service installé) \\
Bureau à distance & Session graphique & Ouverture de session de type~10 \\
Instrumentation à distance & Exécution par l'interface d'administration
Windows & Processus enfant inhabituel \\
Réutilisation d'empreinte & Authentification NTLM avec le haché, sans le
mot de passe & Ouverture de session réseau avec NTLM depuis une machine
inhabituelle \\
\end{tableaucours}
\vspace{-0.8em}
\begin{center}{\small\itshape Tableau~12.2. Techniques de mouvement latéral et traces}\end{center}

% ============================================================================
\section{La Persistance}
\label{sec:c12-persistance}
% ============================================================================

\begin{tableaucours}{L{4.0cm}L{4.4cm}L{3.8cm}}{%
    \tableauHeader{Mécanisme} &
    \tableauHeader{Exemple} &
    \tableauHeader{Trace}
}
Compte créé ou modifié & Compte ajouté aux Admins du domaine & 4720, 4728 \\
Service & Service installé au démarrage & 7045 \\
Tâche planifiée & Tâche lancée chaque nuit & 4698 \\
Clé de démarrage & Programme lancé à l'ouverture de session & Surveillance
du registre \\
Accès distant tiers & Outil de prise de main installé & Inventaire
logiciel \\
\end{tableaucours}
\vspace{-0.8em}
\begin{center}{\small\itshape Tableau~12.3. Mécanismes de persistance}\end{center}

% ============================================================================
\section{Les Limites du Testeur}
\label{sec:c12-limites}
% ============================================================================

\begin{boxalert}
\textbf{Violation typique.} Installer une persistance « pour montrer
qu'on pourrait revenir ».

\smallskip
\textbf{Règle~:} le testeur décrit la persistance possible, il ne l'installe
pas, sauf exercice d'équipe rouge expressément convenu et tracé
(Chapitre~\ref{chap:rouge}). Tout compte, fichier ou service créé est listé
et supprimé avant la fin de la mission, et la suppression est vérifiée.
\end{boxalert}

\begin{boxprocedure}
\textbf{Le nettoyage de fin de mission}
\begin{enumerate}
    \item Relire le journal de mission~: tout ce qui a été créé ou modifié.
    \item Supprimer comptes, fichiers, services, tâches créés.
    \item Restaurer les configurations modifiées.
    \item Faire vérifier la suppression par le client.
    \item Consigner le nettoyage dans le rapport.
\end{enumerate}
\end{boxprocedure}

% ============================================================================
\section{Ce que Voient les Journaux}
\label{sec:c12-journaux}
% ============================================================================

\begin{boxscenario}{Opération LOGONE. Ce que le SOC a vu, et pas vu}
\textbf{Constatations.} Le SOC de nuit a reçu les journaux des contrôleurs
de domaine et des serveurs de Garoua. Il a signalé le balayage. Il n'a pas
corrélé la connexion du compte de la facturation, inhabituelle, à un poste
ordinaire, ni sa nouvelle appartenance au groupe visé.

\medskip
\textbf{Analyse} (confiance moyenne). Les événements utiles étaient
présents~; il manquait les règles. Après l'Opération GAROUA, c'est la
seconde fois que les journaux contiennent l'attaque sans que personne ne la
lise.
\end{boxscenario}

\begin{boxobjet}{T-08 · Comptes d'administration utilisés sur des postes ordinaires}
\textbf{Description~:} des comptes membres de groupes à privilèges ouvrent
des sessions sur des postes de travail. \textbf{Preuve~:} sessions
relevées sur 7~postes de Garoua (26~octobre). \textbf{Gravité~:} moyenne.
\textbf{Impact~:} exposition de secrets à privilèges en cas de
compromission d'un poste. \textbf{Recommandation~:} appliquer le modèle par
niveaux, administrer par le bastion. \textbf{Détection~:} ouverture de
session d'un compte à privilèges sur un poste du niveau inférieur.
\end{boxobjet}

\begin{boxmetier}{Analyste SOC et équipe rouge}
\textbf{Écrire une règle par technique jouée.} Après chaque test, l'analyste
reprend les techniques ATT\&CK jouées et vérifie, une par une, qu'une règle
les aurait vues~; il écrit celles qui manquent.
\end{boxmetier}

\separateur

\begin{boxresume}
\textbf{Ce qu'il faut retenir de ce chapitre~:}
\begin{itemize}
    \item \textbf{Comprendre, s'étendre, durer, atteindre, agir}~: les
          objectifs après l'élévation.
    \item \textbf{Le mouvement latéral} utilise des identifiants valides~;
          il laisse des traces d'ouverture de session et de services.
    \item \textbf{La persistance}~: comptes, services, tâches, clés, outils
          d'accès.
    \item \textbf{Le testeur décrit, n'installe pas}~; il nettoie et fait
          vérifier.
    \item \textbf{Les journaux contiennent l'attaque}~; il faut les règles
          pour la lire.
\end{itemize}
\end{boxresume}

\begin{boxinfo}
\textbf{Points de contrôle de l'élève.} Je sais~: relier les objectifs aux
tactiques ATT\&CK~; associer une technique latérale à sa trace~; citer les
mécanismes de persistance et leurs événements~; dérouler le nettoyage de fin
de mission.
\end{boxinfo}

\bigskip

\begin{boxexo}[title={\ding{45}\quad Exercice~12.1. Tactiques \niveaubase}]
Rangez dans les tactiques d'ATT\&CK les actions de l'attaquant de mars entre
le 3 et le 5~mars.
\end{boxexo}

\begin{boxexo}[title={\ding{45}\quad Exercice~12.2. Événements \niveaubase}]
Quel événement Windows signale~: la création d'un compte~; l'installation
d'un service~; la création d'une tâche planifiée~?
\end{boxexo}

\begin{boxexo}[title={\ding{45}\quad Exercice~12.3. Nettoyage \niveaubase}]
Pourquoi faire vérifier le nettoyage par le client~?
\end{boxexo}

\begin{boxexo}[title={\ding{45}\quad Exercice~12.4. Réutilisation d'empreinte \niveaumoyen}]
Pourquoi la réutilisation d'empreinte rend-elle inutile la longueur du mot
de passe~? Quelle mesure la bloque~?
\end{boxexo}

\begin{boxexo}[title={\ding{45}\quad Exercice~12.5. Une règle \niveaumoyen}]
Écrivez en français la règle qui aurait vu la connexion du compte de la
facturation sur un poste ordinaire.
\end{boxexo}

\begin{boxexo}[title={\ding{45}\quad Exercice~12.6. Persistance tracée \niveaumoyen}]
Dans un exercice d'équipe rouge, une persistance est convenue. Quelles
règles de traçabilité imposez-vous~?
\end{boxexo}

\begin{boxexo}[title={\ding{45}\quad Exercice~12.7. T-08 \niveaumoyen}]
Expliquez le lien entre T-08 et le modèle d'administration par niveaux du
Chapitre~\ref{chap:ad}.
\end{boxexo}

\begin{boxexo}[title={\ding{45}\quad Exercice~12.8. GAROUA rejoué \niveaumoyen}]
Pour chaque mécanisme de persistance de l'attaquant de mars, donnez
l'événement qui l'aurait signalé.
\end{boxexo}

\begin{boxexo}[title={\ding{45}\quad Exercice~12.9. Couverture \niveauexpert}]
Construisez la matrice « technique ATT\&CK jouée / règle de détection /
détecté ou non » de l'Opération LOGONE.
\end{boxexo}

\begin{boxexo}[title={\ding{45}\quad Exercice~12.10. Chasse \niveauexpert}]
Proposez une recherche proactive de persistance à mener sur le parc de
LiZbiZ après l'Opération LOGONE.
\end{boxexo}

\bigskip

\begin{boxinfo}
\textbf{Pour aller plus loin.}
\begin{itemize}
    \item MITRE ATT\&CK~\cite{mitreattack}.
\end{itemize}
\end{boxinfo}

\begin{boxinfo}
\textbf{Transition.} Presque tous les chemins passent par un mot de passe~:
le Chapitre~\ref{chap:mdp} traite de leurs attaques et de leurs défenses.
\end{boxinfo}
